\chapter{Optimasi: Turun Gunung dalam Kabut}
\label{ch:optimasi}

Optimizer sudah saya pakai sejak tugas deep learning di semester pertama, tetapi teorinya baru
dibahas secara sistematis di semester terakhir, di kuliah optimasi numerik: mengapa gradient
descent bisa lambat, kapan ia pasti konvergen, dan apa yang dibeli metode Newton dengan biaya
yang mahal. Bab ini menggabungkan kedua sisi itu.

Bab sebelumnya berakhir dengan perintah ``minimalkan risiko empiris''. Sekarang kita membahas
cara melakukannya. Hampir semua model di buku ini, dari regresi linear sampai PINN, dilatih
dengan varian dari satu gagasan: ikuti arah turunan.

Bayangkan Anda berdiri di lereng gunung yang tertutup kabut tebal. Anda ingin turun ke lembah,
tetapi hanya bisa melihat tanah di sekitar kaki. Strategi paling masuk akal adalah merasakan ke
arah mana tanah paling curam menurun, melangkah ke sana, lalu mengulanginya. Inilah gradient
descent. Persoalannya ada pada panjang langkah. Langkah yang terlalu kecil membuat Anda baru
sampai besok pagi. Langkah yang terlalu besar membuat Anda melompati lembah dan mendarat di
lereng seberang, bahkan lebih tinggi dari tempat semula.

\section{Masalah dan syarat optimalitas}

Kita mencari $\vx^*=\argmin_{\vx\in\R^N}f(\vx)$ untuk fungsi $f$ yang cukup mulus. Dua syarat
klasik membantu mengenali minimum. Kalau $\vx^*$ minimum lokal, gradiennya pasti nol,
$\nabla f(\vx^*)=\bm0$. Sebaliknya, kalau gradiennya nol dan Hessian $\nabla^2f(\vx^*)$ definit
positif, maka $\vx^*$ pasti minimum lokal yang tegas. Titik dengan gradien nol disebut
\istilah{stationary point}, dan ia bisa berupa minimum, maksimum, atau \istilah{saddle point}. Di
dimensi tinggi, saddle point jauh lebih banyak daripada minimum, karena untuk menjadi minimum
semua nilai eigen Hessian harus positif sekaligus.

Fungsi $f$ disebut \istilah{convex} kalau ruas garis yang menghubungkan dua titik di grafiknya
selalu berada di atas grafik itu, yaitu $f(\lambda\vx+(1-\lambda)\vy)\le\lambda f(\vx)+(1-\lambda)f(\vy)$
untuk setiap $\lambda\in[0,1]$. Untuk fungsi konveks, setiap minimum lokal adalah minimum global.
Regresi linear dengan loss kuadrat, regresi logistik, dan lasso semuanya konveks. Jaringan saraf
tidak. Meski begitu, teori konveks tetap berguna, karena di sekitar sebuah minimum hampir semua
fungsi mulus terlihat seperti mangkuk kuadrat.

Fungsi uji klasik yang tidak konveks adalah fungsi Rosenbrock \citep{rosenbrock1960},
$f(x,y)=(1-x)^2+100(y-x^2)^2$. Minimumnya di $(1,1)$, tersembunyi di dasar lembah sempit
berbentuk pisang. Menemukan lembahnya mudah. Yang sulit adalah berjalan menyusuri lembah itu
sampai ke titik terendahnya.

\section{Gradient descent}

Aturan dasarnya adalah
\begin{equation}
\vx_{k+1}=\vx_k-\alpha_k\,\nabla f(\vx_k),
\end{equation}
dengan $\alpha_k>0$ panjang langkah atau \istilah{learning rate}. Arah negatif gradien dipilih
karena ekspansi Taylor orde satu, $f(\vx+\bm h)\approx f(\vx)+\nabla f(\vx)^\top\bm h$, paling
cepat turun kalau $\bm h$ searah $-\nabla f(\vx)$.

Kapan langkah-langkah ini dijamin menuju minimum? Anggap gradien $f$ bersifat Lipschitz dengan
konstanta $L$, artinya $\|\nabla f(\vx)-\nabla f(\vy)\|\le L\|\vx-\vy\|$. Konstanta $L$ mengukur
seberapa cepat kemiringan bisa berubah. Dengan langkah tetap $\alpha=1/L$, setiap iterasi
menjamin
\[
f(\vx_{k+1})\le f(\vx_k)-\frac{1}{2L}\|\nabla f(\vx_k)\|^2,
\]
jadi selama gradien belum nol, nilai fungsi pasti turun.

Ketaksamaan ini, yang di naskah kuliah disebut \istilah{descent lemma}, bisa diturunkan dalam dua baris.
Gradien yang Lipschitz berarti fungsi tidak bisa melengkung ke atas lebih cepat dari sebuah parabola
dengan kelengkungan $L$, sehingga untuk setiap $\bm h$,
\[
f(\vx+\bm h)\le f(\vx)+\nabla f(\vx)^\top\bm h+\frac L2\|\bm h\|^2.
\]
Ruas kanan adalah parabola yang selalu berada di atas fungsi. Masukkan langkah gradient descent,
$\bm h=-\alpha\nabla f(\vx)$, sehingga ruas kanan menjadi
$f(\vx)-\alpha\big(1-\tfrac{L\alpha}2\big)\|\nabla f(\vx)\|^2$. Penurunan terbesar didapat dengan
memaksimalkan $\alpha(1-L\alpha/2)$, yang terjadi pada $\alpha=1/L$, dan nilainya $1/(2L)$. Dari sini juga
terlihat batas atas langkah: kalau $\alpha\ge2/L$, faktor $1-L\alpha/2$ menjadi nol atau negatif, dan
jaminan penurunannya hilang.

Kalau $f$ juga \istilah{strongly convex}
dengan konstanta $\mu$, galat menyusut secara geometris dengan laju yang ditentukan oleh
\istilah{condition number} $\kappa=L/\mu$.

Contoh kecil memperlihatkan arti condition number. Ambil $f(x_1,x_2)=\tfrac12(x_1^2+100x_2^2)$.
Gradiennya $(x_1,\,100x_2)$, sehingga $L=100$, $\mu=1$, dan $\kappa=100$. Dengan langkah tetap
$\alpha$, kedua koordinat berkembang terpisah: $x_1$ dikalikan $(1-\alpha)$ dan $x_2$ dikalikan
$(1-100\alpha)$ di setiap langkah. Supaya $x_2$ tidak meledak, kita butuh $|1-100\alpha|<1$, yang
berarti $\alpha<0{,}02$. Tetapi dengan langkah sekecil itu, $x_1$ hanya menyusut dengan faktor
sekitar $0{,}98$. Bahkan dengan langkah terbaik, $\alpha=2/(L+\mu)$, faktor terburuknya
$(\kappa-1)/(\kappa+1)=99/101$. Untuk mengecilkan galat sejuta kali, dibutuhkan kira-kira
$\ln(10^{-6})/\ln(99/101)\approx690$ langkah. Metode Newton menyelesaikan masalah kuadrat yang
sama dalam satu langkah.

Itulah inti masalah optimasi dalam praktik. Gradient descent lambat bukan karena fungsinya rumit,
melainkan karena skala arah-arahnya berbeda jauh. Pada lembah yang lonjong, lintasannya
berzig-zag menyeberangi lembah alih-alih berjalan menyusurinya (\cref{fig:zigzag}).

\begin{figure}[ht]
\centering
\begin{tikzpicture}[scale=0.9]
  \foreach \r in {0.5,1.0,1.5,2.0,2.5}{
    \draw[abu!60] (0,0) ellipse ({2.6*\r} and {0.55*\r});
  }
  \fill (0,0) circle (1.5pt);
  \draw[oranye,thick,->] (-5.8,1.0) -- (-5.0,-1.0) -- (-4.2,0.85) -- (-3.5,-0.72) -- (-2.9,0.6) -- (-2.4,-0.5) -- (-2.0,0.4) -- (-1.6,-0.3);
  \draw[biru,thick,->] (-5.8,1.0) .. controls (-4.6,-0.4) and (-2.6,-0.2) .. (-0.3,0.0);
  \node[oranye,font=\footnotesize,anchor=west] at (-1.4,-0.9) {gradient descent};
  \node[biru,font=\footnotesize,anchor=west] at (0.2,0.55) {dengan momentum};
\end{tikzpicture}
\caption{Pada lembah lonjong, gradient descent berzig-zag. Momentum meredam gerak bolak-balik
dan mempercepat gerak sepanjang lembah. Gambar skematik.}
\label{fig:zigzag}
\end{figure}

\section{Line search}

Daripada menebak $\alpha$ sekali untuk selamanya, kita bisa mencarinya di setiap iterasi.
\istilah{Line search} mencoba beberapa panjang langkah sepanjang arah turun $\bm d_k$ dan menerima
yang pertama dianggap cukup baik. Ukuran ``cukup baik'' yang paling umum adalah syarat
Armijo \citep{armijo1966},
\[
f(\vx_k+\alpha\bm d_k)\le f(\vx_k)+c\,\alpha\,\nabla f(\vx_k)^\top\bm d_k,\qquad 0<c<1,
\]
yang meminta agar penurunan nyata setidaknya sebagian kecil dari penurunan yang dijanjikan oleh
garis singgung. Prosedur \istilah{backtracking} memulai dari $\alpha=1$ lalu terus mengalikannya
dengan sebuah faktor $\beta\in(0,1)$ sampai syarat itu terpenuhi. Di deep learning, line search
jarang dipakai, karena setiap evaluasi $f$ berarti melewati seluruh data. Gagasannya muncul
kembali dalam L-BFGS.

\section{Metode Newton dan kerabatnya}

Metode Newton memanfaatkan kelengkungan. Ia meminimalkan model kuadrat lokal
$f(\vx_k)+\nabla f^\top\bm h+\tfrac12\bm h^\top\nabla^2f\,\bm h$, yang menghasilkan
\begin{equation}
\vx_{k+1}=\vx_k-\big[\nabla^2 f(\vx_k)\big]^{-1}\nabla f(\vx_k).
\end{equation}
Di dekat minimum yang Hessian-nya definit positif, Newton konvergen secara kuadratik: jumlah
digit yang benar kira-kira berlipat dua di setiap iterasi. Ada tiga hal yang membatasinya.
Hessian berukuran $N\times N$, sehingga untuk jaringan dengan ratusan juta parameter,
menyimpannya saja sudah mustahil. Jauh dari minimum, Hessian bisa tidak definit dan langkah
Newton justru menanjak. Kuliah menyebut perbaikannya \istilah{Hessian modification}, yaitu
menambahkan $\lambda\bm I$ sampai matriksnya definit positif. Terakhir, Newton tertarik ke titik
mana pun yang gradiennya nol, termasuk saddle point.

Metode \istilah{quasi-Newton} seperti BFGS, dan versi hemat memorinya L-BFGS, membangun
aproksimasi Hessian dari selisih gradien beberapa iterasi terakhir \citep{nocedal2006}. L-BFGS
populer untuk masalah deterministik berukuran sedang, termasuk tahap akhir pelatihan PINN.

\section{Stochastic gradient descent}

Dalam machine learning, $f$ adalah rata-rata loss atas $N$ contoh, jadi menghitung gradien penuh
berarti melewati seluruh data. \istilah{Stochastic gradient descent} (SGD) menggantinya dengan
gradien dari sebuah \istilah{mini-batch} acak berukuran $B$:
\begin{equation}
\vw_{k+1}=\vw_k-\alpha_k\,\frac1B\sum_{n\in\mathcal B_k}\nabla_{\vw}L\big(y_n,f(\vx_n;\vw_k)\big).
\end{equation}
Gradien mini-batch adalah estimator tak bias dari gradien penuh, tetapi berisik. Teorinya
berakar pada \istilah{stochastic approximation} \citep{robbins1951}: agar iterasi konvergen,
panjang langkah harus mengecil dengan syarat $\sum_k\alpha_k=\infty$ dan
$\sum_k\alpha_k^2<\infty$. Syarat pertama menjamin kita masih bisa berjalan sejauh apa pun,
syarat kedua menjamin noise akhirnya teredam.

Untuk mengetahui harga rata-rata beras di seluruh kota, Anda bisa mendatangi semua warung, atau
bertanya ke sepuluh warung yang dipilih acak. Jawaban sepuluh warung tidak tepat, tetapi didapat
jauh lebih cepat. Kalau survei kecil itu diulang terus dan tebakan disesuaikan sedikit demi
sedikit, Anda sampai ke angka yang benar lebih cepat daripada menunggu satu survei besar. SGD
bekerja dengan cara yang sama.

Noise SGD ternyata bukan sekadar kekurangan. Ia membuat iterasi terus bergoyang dan sulit
bertahan di minimum yang sempit. \citet{hochreiter1997flat} berargumen bahwa minimum yang datar
generalisasi lebih baik, karena perubahan kecil pada bobot tidak banyak mengubah loss, sehingga
bobotnya bisa dideskripsikan dengan presisi lebih rendah dan modelnya efektif lebih sederhana.
\citet{keskar2017} mengamati bahwa batch yang besar cenderung berakhir di minimum tajam
dengan generalisasi lebih buruk. Kuliah deep learning menyebut efek ini \istilah{self-regularization}
oleh SGD.

\section{Momentum dan metode adaptif}

Metode \istilah{heavy ball} dari \citet{polyak1964} menambahkan ingatan tentang arah sebelumnya,
\begin{equation}
\bm v_{k+1}=\beta\,\bm v_k-\alpha\,\nabla f(\vx_k),\qquad \vx_{k+1}=\vx_k+\bm v_{k+1},
\end{equation}
dengan $\beta$ biasanya sekitar $0{,}9$. Komponen gradien yang berganti tanda setiap langkah saling
meniadakan di dalam $\bm v$, sedangkan komponen yang konsisten menumpuk. Gradient descent biasa
mirip orang yang berhenti total di setiap langkah lalu melihat lagi ke mana arah turun. Momentum
mirip bola yang menggelinding: ia membawa kecepatannya, tidak langsung berbalik karena satu
gundukan kecil, dan melaju lebih cepat di lereng yang panjang. Zig-zag di \cref{fig:zigzag}
teredam dengan cara ini.

Masalah condition number juga bisa diserang dengan memberi setiap parameter laju belajarnya
sendiri. AdaGrad membagi langkah dengan akar jumlah kuadrat gradien masa lalu, RMSProp memakai
rata-rata bergerak agar pembagi itu tidak terus membesar, dan Adam \citep{kingma2015}
menggabungkan momentum dengan RMSProp:
\begin{equation}
\begin{aligned}
\bm m_k &= \beta_1\bm m_{k-1}+(1-\beta_1)\bm g_k, &
\bm v_k &= \beta_2\bm v_{k-1}+(1-\beta_2)\bm g_k^2,\\
\hat{\bm m}_k &= \frac{\bm m_k}{1-\beta_1^k},\quad \hat{\bm v}_k=\frac{\bm v_k}{1-\beta_2^k}, &
\vw_{k+1} &= \vw_k-\alpha\,\frac{\hat{\bm m}_k}{\sqrt{\hat{\bm v}_k}+\epsilon}.
\end{aligned}
\end{equation}
Semua operasinya per elemen, dengan nilai bawaan $\beta_1=0{,}9$, $\beta_2=0{,}999$, dan
$\epsilon=10^{-8}$. Faktor $1-\beta^k$ disebut \istilah{bias correction}, dan asalnya mudah ditelusuri.
Membuka rekursi $\bm m_k$ dari $\bm m_0=\bm0$ memberi
$\bm m_k=(1-\beta_1)\sum_{j=1}^k\beta_1^{k-j}\bm g_j$. Kalau gradien punya ekspektasi tetap
$\E[\bm g]$, maka
\[
\E[\bm m_k]=(1-\beta_1)\big(1+\beta_1+\cdots+\beta_1^{k-1}\big)\E[\bm g]=(1-\beta_1^k)\,\E[\bm g],
\]
karena deret geometri berhingga $\sum_{j=0}^{k-1}\beta_1^j=(1-\beta_1^k)/(1-\beta_1)$. Jadi $\bm m_k$ bias ke
bawah dengan faktor tepat $1-\beta_1^k$, dan membaginya dengan faktor itu menghapus bias. Pada langkah
pertama dengan $\beta_1=0{,}9$, $\bm m_1=0{,}1\,\bm g_1$, sepuluh kali terlalu kecil, dan koreksi
mengembalikannya ke $\bm g_1$. Untuk $\beta_2=0{,}999$, faktornya pada langkah pertama $0{,}001$, sehingga
tanpa koreksi langkah awal Adam akan besar sekali dan tidak menentu. Setelah beberapa ribu langkah,
kedua faktor mendekati satu dan koreksinya tidak berpengaruh lagi.

Beberapa salah paham tentang Adam cukup umum. Adam tidak selalu lebih baik dari SGD. Untuk
beberapa tugas visi, SGD dengan momentum dan jadwal laju belajar yang baik masih menghasilkan
generalisasi yang lebih baik. Menambahkan suku L2 ke loss lalu memakai Adam juga tidak sama
dengan weight decay, karena suku itu ikut dibagi $\sqrt{\hat{\bm v}}$ sehingga efeknya tidak
seragam antar parameter. AdamW memisahkan weight decay dari langkah adaptif untuk mengatasi hal
ini \citep{loshchilov2019}.

Laju belajar yang tetap pun jarang optimal. Pola yang umum adalah \istilah{warm-up}, yaitu mulai
kecil dan naik selama beberapa ratus langkah, lalu menurun perlahan, misalnya mengikuti kurva
kosinus. Warm-up mencegah langkah awal yang terlalu agresif ketika statistik Adam belum stabil.

\section{Lanskap loss jaringan dalam}

Lanskap loss jaringan dalam ternyata berbeda dari intuisi dimensi rendah. Untuk fungsi acak di
dimensi tinggi, sebagian besar titik kritis dengan loss tinggi adalah saddle point, dan
minimum lokal cenderung punya nilai loss yang mirip satu sama lain. Pada jaringan yang sangat
lebar, ada banyak sekali solusi yang sama-sama mencapai loss latih hampir nol, dan
solusi-solusi itu sering terhubung oleh lembah dengan loss rendah. Pertanyaan pentingnya pun
bergeser, dari ``apakah kita menemukan minimum global?'' menjadi ``minimum mana yang dipilih, dan
apakah ia generalisasi?''.

\section{Optimasi di luar kelas}

Contoh mangkuk lonjong di atas bukan sekadar latihan. Masalah condition number yang sama muncul di
banyak tempat dalam buku ini dengan nama yang berbeda. Filter adaptif LMS dalam pemrosesan sinyal adalah
SGD dengan mini-batch satu sampel, dan kecepatan konvergensinya ditentukan oleh sebaran nilai eigen
matriks autokorelasi input (\cref{ch:sinyal}). Algoritme RLS, yang memperbarui invers matriks itu, adalah
kerabat metode Newton. Faktorisasi matriks dalam sistem rekomendasi dilatih dengan SGD
(\cref{ch:rekomendasi}), dan model molekul dilatih dengan Adam seperti jaringan lain.

Di simulasi, loss sebuah PINN adalah jumlah beberapa suku: residu persamaan diferensial, syarat batas,
syarat awal, dan data. Suku-suku ini punya skala dan kelengkungan yang jauh berbeda, sehingga condition
number masalahnya buruk sekali. \citet{wang2021} menunjukkan bahwa gradien suku residu bisa mendominasi
suku batas, sehingga jaringan memenuhi persamaan di dalam domain tetapi melupakan syarat batasnya. Resep
yang sering dipakai adalah memulai dengan Adam untuk keluar dari daerah awal yang buruk, melanjutkan
dengan L-BFGS untuk turun tajam ke minimum, dan menyeimbangkan bobot suku-suku loss secara adaptif
(\cref{ch:pinn}).

\sumberbab{Teori konvergensi, line search, dan metode Newton mengikuti naskah kuliah Numerical
Optimization for AI. Bagian SGD, momentum, dan optimizer adaptif mengikuti bab metode belajar di
kuliah deep learning dasar. Bacaan lanjut: \citet{nocedal2006} untuk optimasi klasik dan
\citet{bottou2018} untuk optimasi skala besar dalam machine learning.}
